\section{Execution Isolation and Determinism}\label{sec:execution}
Each run is a fresh child process (\texttt{python -s -P -B}, an allowlisted environment, a private working directory), wrapped on macOS in \texttt{sandbox-exec} with a deny-network, write-confined profile; resource limits (CPU, file size, open files) are set in the child and a parent watchdog kills the process group. The platform was one Apple M4 Pro machine (macOS 26.7, arm64) with CPython 3.13.15. Memory limits rely on parent polling; address-space and process-count limits are not enforced on this platform; reads are not confined; a child that detaches into a new session could escape the group kill (untested); \texttt{sandbox-exec} is deprecated. Every record carries a flag that the isolation policy is below the specification minimum, which the specification classes as synthetic or benchmark-only grade. \textbf{Linux, Docker and OS-level isolation were not implemented, not tested and not validated. No claim of safe execution of untrusted code and no cross-platform claim is made.}

\textbf{Determinism controls.} Children run with \texttt{PYTHONHASHSEED=0}, \texttt{TZ=UTC}, \texttt{LC\_ALL=C.UTF-8} and single-threaded numerics. \textbf{Determinism is not complete.} The hash seed is fixed and never varied, so hash-order (set-iteration) sensitivity is untested and a fixed seed removes that variation by construction without showing it harmless; \texttt{datetime.now} is not virtualised, so wall-clock-dependent programs can still yield false witnesses; process ids, thread scheduling and addresses are not controlled; and witness identity is timing-dependent for inputs that sit near a resource limit (Section~\ref{sec:failure}). Experiment E2 (hash-order sensitivity) was deliberately not run.
